% =====================================================================
%  1bat-ud2-5.tex  —  HORA 5: Valor numèric: jerarquia i casos difícils
%  (sense preàmbul: s'inclou des de main.tex)
% =====================================================================
\horatitol{Sessió 5 — Valor numèric: jerarquia i casos difícils}%
{Consolidar el càlcul del valor numèric detectant errors típics, treballant
negatius, fraccions i dues variables.}

\subsection{Idea clau}

A la sessió anterior vam calcular valors numèrics en contextos de copagament.
Avui comencem detectant els \textbf{errors més habituals} i, a continuació,
practiquem els casos que més costen: negatius dins d'un quadrat, fraccions
i expressions amb dues variables.

\begin{idea}
\textbf{Errors típics en el càlcul del valor numèric:}
\begin{itemize}[leftmargin=1.4em, itemsep=3pt]
  \item \textbf{Oblidar elevar al quadrat primer.} Si $P=-3x^2$ i $x=2$,
        el resultat és $-3\per 4 = -12$, \emph{no} $(-6)^2=36$.
  \item \textbf{Perdre el signe d'un terme negatiu.} Si $Q=x^2-5x+2$ i
        $x=3$, el terme $-5x$ val $-15$, no $+15$.
  \item \textbf{Confondre $(-x)^2$ amb $-x^2$.}
        $(-3)^2 = 9$ però $-(3^2) = -9$: el parèntesi ho canvia tot.
  \item \textbf{Errors amb fraccions.} Si $x=\tfrac{1}{2}$, llavors
        $x^2 = \tfrac{1}{4}$, \emph{no} $\tfrac{1}{2}$.
\end{itemize}
\textbf{Recordatori de la jerarquia:} potències → productes i divisions →
sumes i restes. Els parèntesis sempre primer.
\end{idea}

\subsection{Exemples}

\begin{exemple}[caça l'error — tres resolucions del docent]
Observa les tres resolucions següents. Decideix quines són correctes i, en les
errònies, identifica \emph{exactament} en quin pas s'ha comès l'error.

\medskip
\begin{tabular}{@{}p{0.9\linewidth}@{}}
\textbf{Cas A.} $P(x) = 2x^2 - 3x + 1$, \; $x = -2$ \\
\quad Resolució: $P(-2) = 2\per(-2)^2 - 3\per(-2)+1 = 2\per 4+6+1 = 15$. \\[4pt]
\textbf{Cas B.} $Q(x) = x^2 - 4x$, \; $x = 3$ \\
\quad Resolució: $Q(3) = 3^2 - 4\per 3 = 9-12 = -3$. \\[4pt]
\textbf{Cas C.} $R(x) = -x^2 + 5$, \; $x = -1$ \\
\quad Resolució: $R(-1) = -(-1)^2+5 = -(-1)+5 = 1+5 = 6$.
\end{tabular}

\medskip
\textbf{Resposta:}
\begin{itemize}[leftmargin=1.4em, itemsep=2pt]
  \item \textbf{Cas A — Correcte.} $(-2)^2=4$; $-3\per(-2)=+6$: tots els passos
        estan ben fets.
  \item \textbf{Cas B — Correcte.} Jerarquia respectada: primer el quadrat,
        després la multiplicació, finalment la resta.
  \item \textbf{Cas C — Error!} $(-1)^2 = +1$ (no $-1$); el resultat correcte
        és $R(-1)=-1+5=4$. L'error és confondre $(-1)^2$ amb $(-1)$.
\end{itemize}
\end{exemple}

\begin{exemple}[expressió amb dues variables — índex de vulnerabilitat]
Un servei social usa l'expressió
\[
V(r, n) = -\num{0,5}\,r + 10n + 20,
\]
on $r$ és la renda mensual familiar (en centenars d'euros) i $n$ és el nombre
de membres de la família. Calcula $V(6, 4)$ (família de 4 persones amb renda
$600\eur$).
\[
V(6,4) = -\num{0,5}\per 6 + 10\per 4 + 20 = -3 + 40 + 20 = 57.
\]
Com més alt és l'índex, \emph{més} vulnerable és la família.
\end{exemple}

\subsection{Exercicis resolts}

\begin{exresolt}[valor numèric amb negatiu al quadrat]
Calcula el valor numèric de $C(x) = \num{0,05}\,x^2 - 10x + 500$ per a $x = -20$.
(Simula el cas d'una «renda negativa» — un deute — per entendre com actua el quadrat.)
\solucio
Substituïm $x = -20$ i respectem la jerarquia:
\begin{align*}
C(-20) &= \num{0,05}\per(-20)^2 - 10\per(-20) + 500 \\
       &= \num{0,05}\per 400 + 200 + 500 &&\text{(primer la potència: $(-20)^2=400$)} \\
       &= 20 + 200 + 500                  &&\text{(producte: $\num{0,05}\per400=20$)} \\
       &= 720.                            &&\text{(suma final)}
\end{align*}
Observació: $(-20)^2 = 400$, igual que $20^2$; el quadrat \emph{sempre} dóna
un resultat positiu.
\resposta{$C(-20) = 720\eur$}
\end{exresolt}

\begin{exresolt}[valor numèric amb fracció]
Calcula el valor numèric de $B(x) = -20x + 600$ per a $x = \dfrac{3}{2}$
(renda en milers d'euros: $1500\eur$).
\solucio
\begin{align*}
B\!\left(\tfrac{3}{2}\right) &= -20\per\tfrac{3}{2} + 600 \\
                             &= -\tfrac{60}{2} + 600 &&\text{(producte de fracció)} \\
                             &= -30 + 600 \\
                             &= 570.
\end{align*}
\resposta{$B\!\left(\tfrac{3}{2}\right) = 570\eur$ de prestació.}
\end{exresolt}

\subsection{Exercicis per fer}

\begin{exercicis}
\begin{exlist}
  \item Calcula el valor numèric de $P(x) = 3x^2 - 2x + 5$ per a:
        \quad (a) $x = 0$ \quad (b) $x = 1$ \quad (c) $x = -1$.

  \item Una associació rep una subvenció modelitzada per $S(x) = -15x + 450$,
        on $x$ és el nombre d'incompliments detectats en l'auditoria. Calcula
        $S(0)$, $S(10)$ i $S(30)$ i interpreta el que passa quan $x=30$.

  \item Calcula el valor numèric de la fórmula de copagament
        $C(x) = \num{0,05}x^2 - 10x + 500$ per a:
        \quad (a) $x = -10$ \quad (b) $x = \dfrac{1}{2}\per 100 = 50$.
        \newline \textbf{(Compte amb el signe):} a l'apartat (a), fixa't com el
        quadrat elimina el signe negatiu.

  \item L'índex de risc d'exclusió d'un casal de barri es calcula amb
        \[
        E(p, f) = 4p - 3f + 12,
        \]
        on $p$ és el nombre de places lliures i $f$ el nombre de famílies en
        llista d'espera. Calcula $E(2, 5)$ i $E(0, 0)$.

  \item \textbf{(Compte amb la jerarquia)} Calcula el valor numèric de
        $T(x) = -2x^2 + 4x - 1$ per a $x = -3$. Resol-lo dos cops: primer
        aplicant la jerarquia correctament, i després \emph{oblidant} elevar
        al quadrat primer (posant $-2\per(-3)$ en lloc de $-2\per(-3)^2$).
        Compara els dos resultats i explica quin pas és crític.

  \item \textbf{(Autoavaluació)} Marca amb una creu la casella que descriu el
        teu nivell a cada habilitat d'aquesta sessió.
        \begin{center}
        \renewcommand{\arraystretch}{1.4}
        \begin{tabular}{|p{6.2cm}|c|c|c|}
        \hline
        \textbf{Habilitat} & \textbf{Domino} & \textbf{Quasi} & \textbf{Em costa} \\
        \hline
        Substituir un valor negatiu & $\square$ & $\square$ & $\square$ \\
        \hline
        Calcular $(-x)^2$ correctament & $\square$ & $\square$ & $\square$ \\
        \hline
        Substituir una fracció & $\square$ & $\square$ & $\square$ \\
        \hline
        Respectar la jerarquia (potències primer) & $\square$ & $\square$ & $\square$ \\
        \hline
        Calcular amb dues variables & $\square$ & $\square$ & $\square$ \\
        \hline
        \end{tabular}
        \end{center}

\textbf{Atenció a la diversitat:} Si marques «em costa» a les tres primeres
habilitats, comença per l'exercici~5.1 amb $x$ enters positius i posa cada
pas en una línia separada. Si ja domines tot el bloc, fes l'apartat (a)
de l'exercici~5.3 també per a $x = -\dfrac{1}{2}$ \textbf{(Ampliació)}.
\end{exlist}
\end{exercicis}

% ---------------------------------------------------------------------
\fullsolucions{Sessió 5}
\begin{solucions}
\begin{sollist}
  \item (a) $5$ \quad (b) $6$ \quad (c) $10$
  \item $S(0)=450\eur$; \; $S(10)=300\eur$; \; $S(30)=0\eur$:
        amb 30 incompliments la subvenció s'esgota completament.
  \item (a) $C(-10) = \num{0,05}\per 100 + 100 + 500 = 5+100+500 = 605\eur$
        \quad (b) $C(50) = \num{0,05}\per\num{2500} - 500 + 500 = 125\eur$
  \item $E(2,5) = 8-15+12 = 5$ \quad $E(0,0) = 12$
  \item Correcte: $T(-3) = -2\per 9 + 4\per(-3)-1 = -18-12-1 = -31$. \\
        Amb l'error: $-2\per(-3)+4\per(-3)-1 = 6-12-1 = -7$. \\
        El pas crític és elevar $(-3)^2=9$ \emph{abans} de multiplicar per $-2$.
  \item Autoavaluació personal (no hi ha una resposta única).
\end{sollist}
\end{solucions}
